% !TeX root = ../教材题库_全部条目_转写工作稿.tex
% Source: XJ-B1-C05-RV-P209-Y0333
% 第5章 / 章末复习 复习题五 / 教材 p.209 / PDF 第216页
% OCR未校对：公式与文字均可能有误；未完成数学TeX转写；禁止直接发给学生
% 原题及相邻原解截图（核对用）：../题目原图/XJ-B1-C05-RV-P209-Y0333-01.webp
% 本文件只在转写完成后作为题干使用；OCR 条目尚不是可靠题干。
\begin{bankitem}{复习题五 第23题}{第5章 / 章末复习 复习题五 / 教材 p.209 / PDF 第216页}{XJ-B1-C05-RV-P209-Y0333}{OCR未校对：公式与文字均可能有误；未完成数学TeX转写；禁止直接发给学生}
\begin{ocrdraft}
[教材p.209 / PDF第216页]
\par
23．生成于大西洋的强烈热带气旋被称为飓
\par
风．中心风速178、209km/h对应于三级飓风，中
\par
风速210、249km/h对应于四级飓风，中心风速
\par
超过250km/h对应于五级飓风．以下数据是大西
\par
洋海域从1921年到2010年每十年的主要飓风（含
\par
、四、五级）数量．
\par
时间/年
\par
主要飓风数量Y
\par
1921一1930
\par
1
\par
193r一1940
\par
2
\par
16
\par
1941一1950
\par
3
\par
29
\par
1951一1960
\par
4
\par
33
\par
1961．一1970
\par
5
\par
27
\par
1971一1980
\par
6
\par
16
\par
1981一1990
\par
7
\par
16
\par
1991一一2000
\par
8
\par
27
\par
200r一2010
\par
9
\par
33
\par
（D根据上表，试绘制“带平滑线和数据的散点图”；
\par
（2）借助图象，选用一个三角函数模型来近似刻画3，与之间的对应关系，
\par
并求出该拟合模型的函数解析式．
\end{ocrdraft}
\end{bankitem}
